% !PRINTED 212
\weberpage{212}
%<<CONT
\weq{10}{S[F_0(x)] = 0, \quad S[F_1(x)] = 0, \ \ldots \ S[F_{n-2}(x)] = 0.}
If, in accordance with (3), we set
\begin{align*}
F(t, x) &= \frac{f(t)}{t-x} - \frac{1}{n} f'(t) \tag{11}\\
&= t^{n-2} F_0(x) + t^{n-3} F_1(x) + \cdots + t F_{n-3}(x) + F_{n-2}(x),
\end{align*}
then the right-hand side of (8) arises from $F(t, x)$ by the
replacement of $t^k$ by $t_k$.

If we assume $y$ from the outset in the form
\weq{12}{y = t_{n-2} F_0(x) + t_{n-3} F_1(x) + \cdots + t_0 F_{n-2}(x)}
then the equation
\weq{13}{S(y) = 0}
is identically satisfied. What advantage this form of the substitution
affords, the next considerations will show.

\begin{center}
\S.\,69.

\gesperrt{Invariant Property of the Tschirnhausen Transformation.}
\end{center}

We have now to investigate the influence which a linear
transformation, to which we subject the function $f(x)$, has on the
Tschirnhausen transformation.

We make in $f(x)$ the linear substitution
\weq{1}{x = \frac{\alpha \xi + \beta}{\gamma \xi + \delta}, \quad \alpha \delta - \beta \gamma = r,}
whereby we obtain
\weq{2}{\varphi(\xi) = (\gamma \xi + \delta)^n f\left(\frac{\alpha \xi + \beta}{\gamma \xi + \delta}\right),}
so that $\varphi(\xi)$ is an integral rational function of degree $n$, and
$\varphi(\xi) = 0$ is the equation $f(x) = 0$ transformed by (1).

We now derive a function $\Phi(\tau, \xi)$ from $\varphi(\xi)$ in exactly the same
way as we derived $F(t, x)$ from
$f(x)$ in the preceding paragraph, namely
\begin{align*}
\Phi(\tau, \xi) &= \frac{\varphi(\tau)}{\tau - \xi} - \frac{1}{n} \varphi'(\tau) \tag{3}\\
&= \tau^{n-2} \Phi_0(\xi) + \tau^{n-3} \Phi_1(\xi) + \cdots + \tau \Phi_{n-3}(\xi) + \Phi_{n-2}(\xi),
\end{align*}
in which $\tau$, just like $t$, is a variable, and the functions
$\Phi_0(\xi), \Phi_1(\xi) \ldots \Phi_{n-2}(\xi)$, likewise from $\xi$ and the coefficients
%CONT>>
